% SEGMENT OLP-0019-S01
% Part: sets-functions-relations
% Chapter: relations
% Section: operations

\documentclass[../../../include/open-logic-section]{subfiles}

\begin{document}

\olfileid{sfr}{rel}{ops}
\olsection{தொடர்புகள் மீதான செயல்கள்}

தொடர்புகளை மாற்றுவதோ ஒன்றிணைப்பதோ பல நேரங்களில் பயனுள்ளதாக இருக்கிறது.
\olref[sfr][rel][ord]{prop:stricttopartial} இல் தொடர்புகளின்
\emph{சேர்ப்பை} நாம் கருதினோம். இரு தொடர்புகளை ஜோடிகளின் கணங்களாகக்
கருதி எடுக்கும் சேர்ப்பே அது. அதேபோல,
\olref[sfr][rel][ord]{prop:partialtostrict} இல் தொடர்புகளின் சார்பு
வேறுபாட்டைக் கருதினோம். தொடர்புகள் மீது செய்யக்கூடிய வேறு சில செயல்கள்
இங்கே தரப்படுகின்றன.

% SEGMENT OLP-0019-S02
\begin{defn}\ollabel{relationoperations}
$R$, $S$ ஆகியவை தொடர்புகளாகவும், $A$ ஏதேனும் ஒரு கணமாகவும் இருக்கட்டும்.

$R$ இன் \emph{தலைகீழ்த் தொடர்பு}
$R^{-1} = \Setabs{\tuple{y, x}}{\tuple{x,
    y} \in R}$ ஆகும்.

$R$, $S$ ஆகியவற்றின் \emph{சார்புப் பெருக்கம்}
$(R \mid S) =
\{\tuple{x, z} : \exists y(Rxy \land Syz)\}$ ஆகும்.

$R$ ஐ $A$ க்குக் \emph{கட்டுப்படுத்துவது}
$\funrestrictionto{R}{A}= R
\cap A^2$ ஆகும்.

$R$ ஐ $A$ மீது \emph{செயல்படுத்துவதால்} கிடைப்பது
$\funimage{R}{A} = \{y :
(\exists x \in A)Rxy\}$ ஆகும்.
\end{defn}

% SEGMENT OLP-0019-S03
\begin{ex}
$S \subseteq \Int^2$ என்பது $\Int$ மீதான அடுத்துறுப்புத் தொடர்பாக
இருக்கட்டும். அதாவது
$S = \Setabs{\tuple{x, y} \in \Int^2}{x + 1 = y}$; எனவே
$x + 1 = y$ ஆக இருக்கும்போது, அப்போதும் அப்போது மட்டுமே, $Sxy$.

$S^{-1}$ என்பது $\Int$ மீதான முந்துறுப்புத் தொடர்பு; அதாவது
$\Setabs{\tuple{x,y}\in\Int^2}{x -1 =y}$.

$S\mid S$ என்பது
$ \Setabs{\tuple{x,y}\in\Int^2}{x + 2 =y}$.

$\funrestrictionto{S}{\Nat}$ என்பது $\Nat$ மீதான அடுத்துறுப்புத் தொடர்பு.

$\funimage{S}{\{1,2,3\}}$ என்பது $\{2, 3, 4\}$.
\end{ex}

% SEGMENT OLP-0019-S04
\begin{defn}[கடப்பு அடைவு]
$R \subseteq A^2$ ஒரு இருமத் தொடர்பாக இருக்கட்டும்.

$R$ இன் \emph{கடப்பு அடைவு}
$R^+ = \bigcup_{0 < n \in
\Nat} R^n$ ஆகும். இங்கு $R^1 = R$, $R^{n+1} = R^n
\mid R$ என்று மறுநிகழ்வு முறையில் வரையறுக்கிறோம்.

$R$ இன் \emph{தற்சுட்டுக் கடப்பு அடைவு}
$R^* = R^+ \cup
\Id{A}$ ஆகும்.
\end{defn}

% SEGMENT OLP-0019-S05
\begin{ex}
அடுத்துறுப்புத் தொடர்பு $S \subseteq \Int^2$ ஐ எடுத்துக்கொள்வோம்.
$x + 2 =
y$ ஆக இருக்கும்போது, அப்போதும் அப்போது மட்டுமே, $S^2xy$;
$x + 3 = y$ ஆக இருக்கும்போது, அப்போதும் அப்போது மட்டுமே, $S^3xy$;
இவ்வாறே தொடரும். எனவே ஏதேனும் $n \geq 1$ க்கு $x + n = y$ ஆக
இருக்கும்போது, அப்போதும் அப்போது மட்டுமே, $S^+xy$.
வேறு சொற்களில், $x < y$ ஆக இருக்கும்போது, அப்போதும் அப்போது மட்டுமே,
$S^+xy$; மேலும் $x \le
y$ ஆக இருக்கும்போது, அப்போதும் அப்போது மட்டுமே, $S^*xy$.
\end{ex}

% SEGMENT OLP-0019-S06
\begin{prob}
$R$ இன் கடப்பு அடைவு உண்மையிலேயே கடப்புப் பண்பைக் கொண்டுள்ளது எனக் காட்டுக.
\end{prob}

\end{document}
